\subsection{Implementation and validation of the near-vacuum evaporation model}

A near-vacuum evaporation closure based on Wang et al.~\cite{Wang2020}
was implemented as a runtime-selectable model in \texttt{vacuumLaserbeamFoam}.
The implementation evaluates the Knudsen-layer jump state as a function of
Mach number and solves the common-atmosphere transition state at prescribed
surface temperature. The transition temperatures $T_{k0}$ and $T_{k1}$,
corresponding to $\mathrm{Ma}=0.05$ and $\mathrm{Ma}=1$, are determined
automatically. Below the activation temperature, evaporation is disabled;
between the activation temperature and $T_{k1}$ the transition Mach number is
solved iteratively, while the strong-evaporation branch uses $\mathrm{Ma}=1$.
The regular $\mathrm{Ma}\rightarrow0$ endpoint was included explicitly to
avoid a numerical endpoint failure at the chamber-pressure boiling threshold.

For alloys, component mass fractions are converted to molar fractions and the
total saturation pressure and vapor molar mass are evaluated following Wang
et al.'s multicomponent formulation. For 304L stainless steel,
Cr/Ni/Fe=18/8/74 wt\% was used. The mixture saturation-pressure curve was
anchored to $P_e=20.16$ Pa at $T=2009$ K.

The recoil pressure is applied to the VOF interface through the interface
gradient in the momentum equation, while the evaporative heat flux is
localized to the interface in the energy equation. Grey-body radiation is
treated as a separate surface heat-loss term. The chamber pressure enters the
evaporation closure directly; the outer phase remains an incompressible
numerical pseudo-gas and is not intended to represent chamber-scale rarefied
gas dynamics.

\subsection{304L near-vacuum benchmark}

The stationary-laser 304L benchmark was performed at 0.0002 atm
($\approx20.265$ Pa) and 298 K using a 260 W, 100 $\mu$m laser at 1070 nm.
The reference mesh spacing was 4 $\mu$m ($80^3$ cells), and the simulation was
advanced to 140 $\mu$s using 48 MPI ranks. For this validation case, the
inherited Drude/resistivity optical surrogate was replaced by a fixed complex
refractive index for Fe based on Johnson--Christy optical data
\cite{JohnsonChristy1974}. At 1070 nm, the interpolated values were
$n=2.96135$ and $k=4.01327$, corresponding to a normal-incidence single-hit
absorptivity of approximately 0.3725. The deposited laser power at 1 $\mu$s
was 97.12 W, close to the flat-surface Fresnel expectation of 96.85 W.

The formal keyhole depth was extracted from the deepest point of the
atmosphere-connected three-dimensional
$\alpha_{\mathrm{metal}}=0.5$ interface component. The keyhole depth increased
from 32 to 136 $\mu$m in 76.23 $\mu$s. Wang et al. reported approximately
75 $\mu$s for their current evaporation model and approximately 70 $\mu$s
for the corresponding x-ray experiment. An independent centerline extraction
gave 76.02 $\mu$s.

During the 32--136 $\mu$m growth interval, the mean deposited laser power was
212.96 W, the mean evaporative heat loss was 7.76 W, and the mean radiative
loss was 0.0737 W. At the recoil-comparison stage corresponding to 75 $\mu$s
after the 32-$\mu$m-depth reference state, the full connected-surface
pressure-load integral was approximately 3.52 mN, compared with the
approximately 4 mN z-direction recoil force reported by Wang et al. The
time-history peak of this integral was 4.70 mN. The area-weighted 99th
percentile of the surface recoil pressure was approximately 6.29 atm; the
single-face maximum was substantially larger and is retained as a
mesh/interpolation-sensitivity diagnostic rather than used for calibration.

These constitutive, optical, keyhole-growth and recoil-load checks support
freezing the 304L near-vacuum implementation without empirical rescaling of
the evaporation coefficients. Subsequent 0.6 Pa powder-bed simulations use
this closure unchanged.
